%%%%%%%%%%%%%%%%%%%%%%%%%%%%%%%%%%%%%%%%%
% Beamer Presentation
% LaTeX Template
% Version 1.0 (10/11/12)
%
% This template has been downloaded from:
% http://www.LaTeXTemplates.com
%
% License:
% CC BY-NC-SA 3.0 (http://creativecommons.org/licenses/by-nc-sa/3.0/)
%
%%%%%%%%%%%%%%%%%%%%%%%%%%%%%%%%%%%%%%%%%

%----------------------------------------------------------------------------------------
%	PACKAGES AND THEMES
%----------------------------------------------------------------------------------------

\documentclass[9pt,aspectratio=169]{beamer}
\usepackage[utf8]{inputenc}
\usepackage{bm}
\usepackage{bbm}
\usepackage{graphicx}
\usepackage{cmbright}
\usefonttheme{professionalfonts}
%\usefonttheme{serif}
% The Beamer class comes with a number of default slide themes
% which change the colors and layouts of slides. Below this is a list
% of all the themes, uncomment each in turn to see what they look like.

%\usetheme{default}
%\usetheme{AnnArbor}
% \usetheme{Antibes}
%\usetheme{Bergen}
%\usetheme{Berkeley}
%\usetheme{Berlin}
%\usetheme{Boadilla}
% \usetheme{CambridgeUS}
%\usetheme{Copenhagen}
%\usetheme{Darmstadt}
% \usetheme{Dresden}
\usetheme{Frankfurt}
%\usetheme{Goettingen}
% \usetheme{Hannover}
% \usetheme{Ilmenau}
% \usetheme{JuanLesPins}
%\usetheme{Luebeck}
%\usetheme{Madrid}
%\usetheme{Malmoe}
%\usetheme{Marburg}
% \usetheme{Montpellier}
% \usetheme[width=2.5cm]{PaloAlto}
%\usetheme{Pittsburgh}
% \usetheme{Rochester}
%\usetheme{Singapore}
% \usetheme{Szeged}
% \usetheme{Warsaw}
% \setbeamertemplate{title page}[default][colsep=-4bp,rounded=false]
\setbeamertemplate{headline}{}
\defbeamertemplate{footline}{myframe number}
{
  \hspace{1pt}
    \usebeamercolor[structure]{page number in head/foot}%
  \usebeamerfont{page number in head/foot}%
  \raisebox{1.7ex}[0pt][0pt]{% <--- change here
    \insertframenumber\,/\,\inserttotalframenumber\kern1em}%
}
\setbeamertemplate{footline}[myframe number]
\AtBeginSection[]
{
  \begin{frame}
    \frametitle{Outline}
    \begin{minipage}{1.0\linewidth}
      \tableofcontents[currentsection,currentsubsection]
    \end{minipage}
  \end{frame}
}

\AtBeginSubsection[]
{
  \begin{frame}
    \frametitle{Outline}
    \begin{minipage}{1.0\linewidth}
      \tableofcontents[currentsection,currentsubsection]
    \end{minipage}
  \end{frame}
}

% As well as themes, the Beamer class has a number of color themes
% for any slide theme. Uncomment each of these in turn to see how it
% changes the colors of your current slide theme.
%\definecolor{mycolor}{rgb}{.8,.0,.3}
%\setbeamercolor*{palette primary}{use=structure,fg=white,bg=mycolor}
%\usecolortheme{albatross}
%\usecolortheme{beaver}
%\usecolortheme{beetle}
%\usecolortheme{crane}
%\usecolortheme{dolphin}
%\usecolortheme{dove}
%\usecolortheme{fly}
%\usecolortheme{lily}
%\usecolortheme{orchid}
% \usecolortheme{rose}
%\usecolortheme{seagull}
\usecolortheme{seahorse}
%\usecolortheme{whale}
%\usecolortheme{wolverine}
%\usecolortheme[named=mycolor]{structure}
%\setbeamertemplate{footline} % To remove the footer line in all slides uncomment this line
%\setbeamertemplate{footline}[page number] % To replace the footer line in all slides with a simple slide count uncomment this line

%\setbeamertemplate{navigation symbols}{} % To remove the navigation symbols from the bottom of all slides uncomment this line


\usepackage{booktabs} % Allows the use of \toprule, \midrule and \bottomrule in tables

%----------------------------------------------------------------------------------------
%	TITLE PAGE
%----------------------------------------------------------------------------------------

\title[Lecture 3: Unconstrained Optimization]{Lecture 3: Unconstrained Optimization}

\author{Junnan Zhang} % Your name
\institute[Paula and Gregory Chow Institute for Studies in Economics\\
Xiamen University] % Your institution as it will appear on the bottom of every slide, may be shorthand to save space
{
	Paula and Gregory Chow Institute for Studies in Economics \\ Xiamen University  \\ % Your institution for the title page
	\medskip
    \textit{Slides Prepared by Xiaoling Mei}
	%\textit{luciamay@smith.com} % Your email address
}
\date{Fall, 2026} % Date, can be changed to a custom date
\setbeamertemplate{itemize items}[default]
\setbeamertemplate{sidebar}[sidebar right]
% \definecolor{myco}{HTML}{a67678}
\definecolor{myco}{HTML}{8776a6}
\setbeamercolor{itemize item}{fg=myco} % all frames will have red bullets
\setbeamercolor{itemize subitem}{fg=myco} % all frames will have red bullets
\setbeamercolor{block title}{bg=myco}
\setbeamercolor{structure}{fg=myco}
% \setbeamercovered{transparent}
\newcommand{\RR}{\mathbbm R}

\begin{document}



\begin{frame}
\titlepage % Print the title page as the first slide
\end{frame}

\section{First-Order Conditions}

\begin{frame}
\frametitle{Definitions}
Consider a real-valued function of $n$ variables $F$: $U\subset \RR^n
\rightarrow \RR$. A point $\mathbf{x}^*$ is a
\begin{itemize}
    \item \textbf{maximizer} of $F$ on $U$ if
    $F(\mathbf{x}^*)\geq F(\mathbf{x})$ for all $\mathbf{x} \in U$
    \item \textbf{strict maximizer} if $F(\mathbf{x}^*) > F(\mathbf{x})$ for
    all $\mathbf{x}\neq \mathbf{x}^* \in U$
    \item \textbf{local/relative maximizer} if there is a ball $B_r(\mathbf{x}^*)$
    about $\mathbf{x}^*$ such that $F(\mathbf{x}^*) \geq F(\mathbf{x})$ for
    all $\mathbf{x}\in B_r(\mathbf{x}^*)\cap U$
    \item \textbf{strict local maximizer} if there is a ball
    $B_r(\mathbf{x}^*)$ about $\mathbf{x}^*$ such that $F(\mathbf{x}^*) >
    F(\mathbf{x})$ for all $\mathbf{x} \neq \mathbf{x}^* \in
    B_r(\mathbf{x}^*)\cap U$
\end{itemize}
\end{frame}
		

\begin{frame}
	\frametitle{First Order Conditions}
	\begin{block}{\textbf{Theorem 17.1}}	
        Let $F$: $U\subset \RR^n \rightarrow \RR^1$ be a $C^1$ function. If
        $\mathbf{x}^*$ is a local max or min of $F$ in $U$ and if
        $\mathbf{x}^*$ is an interior point of $U$, then
        $$\frac{\partial F}{\partial x_i}(\mathbf{x}^*) = 0\ \ \ \ \ \text{for}\ i =
        1,\cdots, n. $$
	\end{block}
	Proof: for each $x_i$, $x_i^*$ is an interior maximizer of $x_i \mapsto
    F(x_1^*, \ldots, x_{i-1}^*, x_i, x_{i+1}^*, \ldots, x_n^*)$. Use Theorem 3.3.
\end{frame}	


\section{Quadratic Forms (Chapter 16)}
\begin{frame}
	\frametitle{Quadratic Forms}
    \begin{block}{\bfseries Definition}
        \textbf{Definition:} A quadratic form on $\RR^k$ is a
        real-valued function of the form
        $$Q(x_1,\cdots,x_k) = \sum_{i\leq j}a_{ij}x_ix_j$$
        in which each term is a monomial of degree two.
    \end{block}
    A quadratic form can be written in matrix form:
    \begin{align}
        \mathbf{x}^T\! A\, \mathbf{x} = (x_1\ x_2\ \ \cdots\ \
        x_k)\left({\begin{array}{cccc}
                  a_{11} & \frac{1}{2}a_{12} & \cdots & \frac{1}{2}a_{1k}\\
                  \frac{1}{2}a_{12} & a_{22} & \cdots & \frac{1}{2}a_{2k}\\
                  \vdots & \vdots & \ddots & \vdots \\
                  \frac{1}{2}a_{1k} & \frac{1}{2}a_{2k} & \cdots & a_{kk}\\
              \end{array}}\right)\left({\begin{array}{c}
                  x_1 \\
                  x_2 \\
                  \vdots \\
                  x_k \\
              \end{array}}\right) \nonumber
    \end{align}
    where $A$ is a symmetric matrix.
\end{frame}


\begin{frame}
	\frametitle{Definiteness}
    A quadratic form (or a symmetric matrix $A$) is
	\begin{itemize}
		\item positive definite: if $\mathbf{x}^T\!A\,\mathbf{x}>0$ for all $\mathbf{x}\neq0$ in $\RR^n$. For example: $Q(x_1,x_2) = x_1^2+x_2^2$
		 \item negative definite: if $\mathbf{x}^T\!A\,\mathbf{x}<0$ for all $\mathbf{x}\neq0$ in $\RR^n$. For example: $Q(x_1,x_2) = -x_1^2-x_2^2$
		 \item positive semi-definite: if $\mathbf{x}^T\!A\,\mathbf{x}\geq0$ for all $\mathbf{x}\neq0$ in $\RR^n$. For example: $Q(x_1,x_2) = x_1^2+2x_1x_2+x_2^2$
		 \item negative semi-definite: if $\mathbf{x}^T\!A\,\mathbf{x}\leq0$ for all $\mathbf{x}\neq0$ in $\RR^n$. For example: $Q(x_1,x_2) = -x_1^2-2x_1x_2-x_2^2$
		 \item indefinite: if  $\mathbf{x}^T\!A\,\mathbf{x} > 0$ for some $\mathbf{x} $ in $\RR^n$ and $< 0$ for some other $\mathbf{x}$ in $\RR^n$. \\
		 for example: $Q(x_1,x_2) = x_1^2-x_2^2$
	\end{itemize}
\end{frame}


\begin{frame}
	\frametitle{Principal Minors}
    \begin{block}{\bfseries Definition}
        Let $A$ be an $n\times n$ matrix. A $k\times k$ submatrix of $A$
        formed by deleting $n-k$ columns, say columns $i_1,i_2,\cdots,
        i_{n-k}$ and the same rows $i_1,i_2,\cdots, i_{n-k}$ from $A$ is
        called a $k$th order principal submatrix of $A$. The determinant of
        a $k\times k$ principal submatrix is called a $k$th order principal
        minor of A.
    \end{block}
    \begin{block}{\bfseries Definition}
        Let $A$ be an $n\times n$ matrix. The $k$th order principal submatrix
        of $A$ obtained by deleting the last $n-k$ rows and the last $n-k$
        columns from $A$ is called $k$th order leading principal submatrix of
        $A$. Its determinant is called $k$th order \emph{leading principal}
        minor of $A$.
    \end{block}
    For example, for a $3\times 3$ matrix, there are three leading principal
    minors:
    \begin{equation*}
        |a_{11}|,\,
        \begin{vmatrix}
            a_{11} & a_{12} \\
            a_{21} & a_{22}
        \end{vmatrix},\,
        \begin{vmatrix}
            a_{11} & a_{12} & a_{13} \\
            a_{21} & a_{22} & a_{23} \\
            a_{31} & a_{32} & a_{33}
		\end{vmatrix}
    \end{equation*}
\end{frame}


\begin{frame}[allowframebreaks]
	\frametitle{Characterization of Definiteness}
	\begin{block}{\textbf{Theorem 16.1}}
        Let A be an $n\times n$ symmetric matrix. Then,
        \begin{enumerate}
            \item[(a)] $A$ is positive definite \emph{if and only if} all its $n$
            leading principal minors are (strictly) positive
            \item[(b)] $A$ is negative definite \emph{if and only if} its $n$
            leading principal minors alternate in sign as follows: $|A_1|<0$,
            $|A_2|>0$, $|A_3|<0$, etc. That is, the $k$th order leading principal
            minor should have the same sign as $(-1)^k$.
            \item[(c)] If some $k$th order leading principal minor of $A$ is
            nonzero but does not fit either of the above two sign patterns,
            then $A$ is indefinite. This case occurs when $A$ has a negative
            $k$th order leading principal minor for an \emph{even} integer
            $k$; or when $A$ has a negative $k$th order leading principal
            minor and a positive $\ell$th order leading principal minor for
            two distinct \emph{odd} integers $k$ and $\ell$.
        \end{enumerate}
    \end{block}
    This test may fail for a symmetric matrix when some leading principal
    minor is 0 while the nonzero ones fit the pattern in either (a) or (b).

    \framebreak
	\begin{block}{\textbf{Theorem 16.2}}
		Let $A$ be an $n\times n$ symmetric	matrix. Then
		\begin{itemize}
			\item  $A$ is positive semidefinite \emph{if and only if}
			\textbf{every} principal minor of A is $\geq0$
			\item 	$A$ is negative semidefinite \emph{if and only if} \textbf{every} principal minor of odd order is $\leq 0$ and every
			principal minor of even order is $\geq0$.
		\end{itemize}
	\end{block}
\end{frame}

\begin{frame}
	\frametitle{Definiteness of Diagonal Matrices}
\begin{itemize}
	\item The simplest $n\times n$ symmetric matrix is the diagonal matrix,
    which corresponds to the simplest quadratic form $a_1x_1^2+\cdots
    a_nx_n^2$
	\item It is positive definite $\iff$ all the $a_i$'s are positive and
    negative definite $\iff$ all the $a_i$'s are negative
	\item It is positive semidefinite $\iff$ all $a_i \geq 0$ and negative
    semidefinite $\iff$ all $a_i \leq 0$
	\item If there are two $a_i$'s of opposite signs, it will be indefinite
	\item If some $a_j = 0$, it is not definite
\end{itemize}
\end{frame}

\begin{frame}
\frametitle{Definiteness and Optimality}
\begin{itemize}
	\item Fact: determining the definiteness of a quadratic form $Q$ is
    equivalent to determining whether $\mathbf{x} = \mathbf{0}$ is a max, a
    min, or neither for the real-valued function $Q$.
    \item Figures 16.2--16.6
\end{itemize}
\end{frame}

\section{Second-Order and Global Conditions}

\begin{frame}
	\frametitle{Second Order Conditions}
	\begin{itemize}	
        \item A point $x^*$ is a \textbf{critical point} if $DF(\mathbf{x}^*)
        = 0$
		\item To determine if a critical point is a max or min, we need to
        use a condition on the second derivatives of $F$
		\item \textbf{Hessian} of $F$: 
			$$\left({\begin{array}{ccc} 
				\frac{\partial^2 F}{\partial x_1^2} & \cdots &\frac{\partial^2 F}{\partial x_n\partial x_1}\\
				\vdots & \ddots & \vdots\\
				\frac{\partial^2 F}{\partial x_1\partial x_n} & \cdots & \frac{\partial^2 F}{\partial x_n^2}\\			
				\end{array}}\right)$$
	\end{itemize}
\end{frame}

\begin{frame}
	\frametitle{Second Order Conditions}
	\begin{block}{\textbf{Theorem 17.2 (Sufficient Conditions)}}	
		Let $F$: $U\subset \RR^n \rightarrow \RR^1$ be a $C^2$ function whose
        domain is an open set $U$ in $\RR^n$. Suppose $\mathbf{x}^*$ is a
        critical point of $F$:
		\begin{enumerate}
            \item[(a)] if the Hessian $D^2F(\mathbf{x^*})$ is a negative
            definite symmetric matrix, then $\mathbf{x^*}$ is a strict local
            max of $F$;
            \item[(b)] if the Hessian $D^2F(\mathbf{x^*})$ is a positive
            definite symmetric matrix, then $\mathbf{x^*}$ is a strict local
            min of $F$;
            \item[(c)] if the Hessian $D^2F(\mathbf{x^*})$ is indefinite,
            then $\mathbf{x^*}$ is neither a local max nor a local min of
            $F$.
		\end{enumerate} 
	\end{block}
\end{frame}



\begin{frame}
	\frametitle{Second Order Conditions}
	\begin{itemize}%{\textbf{Second order Conditions}}	
        \item Equivalent statements based on analytical characterization of
        positive definite and negative definite matrices: Theorem 17.3-17.5
        \item A critical point of $F$ for which the Hessian
        $D^2F(\mathbf{x^*})$ is indefinite is called a \textbf{saddle point}
        \item A saddle point is a min of $F$ in some directions and a max in
        other directions
        \item Example: $F(x_1,x_2) = x_1^2-x_2^2$
	\end{itemize}
\end{frame}


\begin{frame}
	\frametitle{Necessary Conditions}
	\begin{block}{\textbf{Theorem 17.6}}	
		Let $F$: $U\subset \RR^n \rightarrow \RR^1$ be a $C^2$ function. Suppose $\mathbf{x}^*$ is an interior point of $U$ and a local max/min of $F$. Then  
		$$DF(\mathbf{x^*}) = \mathbf{0}$$ and $D^2F(\mathbf{x^*})$ is negative/positive \textbf{semi}definite. 
	\end{block}
\end{frame}


\begin{frame}
	\frametitle{Necessary Conditions}
	\begin{block}{\textbf{Theorem 17.7}}	
		Let $F$: $U\subset \RR^n \rightarrow \RR^1$ be a $C^2$ function of
        $n$ variables. Suppose $\mathbf{x}^*$ is an interior point of $U$ and
        a local max/min of $F$.
        \begin{enumerate}
            \item[(a)] if $\mathbf{x^*}$ is a local min of $F$, then
            $\frac{\partial F}{\partial x_i}(\mathbf{x^*}) = 0$ for
            $i=1,\cdots,n$ and all the principal minors of the Hessian
            $D^2F(\mathbf{x^*})$ are $\geq 0$.
            \item[(b)] if $\mathbf{x^*}$ is a local max of $F$, then
            $\frac{\partial F}{\partial x_i}(\mathbf{x^*}) = 0$ for
            $i=1,\cdots,n$ and all the principal minors of the Hessian
            $D^2F(\mathbf{x^*})$ of odd order are $\leq 0$ and all the
            principal minors of the Hessian $D^2F(\mathbf{x^*})$ of even
            order are $\geq 0$.
        \end{enumerate}
	\end{block}
\end{frame}

\begin{frame}
	\frametitle{Global Maximum and Minimum}
	\begin{block}{\textbf{Theorem 17.8}}
		Let $F$: $U \rightarrow \RR^1$ be a $C^2$ function whose domain is a convex open subset $U$ of $\RR^n$. 
		\begin{enumerate}
			\item[(a)] The following three conditions are equivalent: 
			\begin{enumerate}
				\item[(i)] $F$ is a concave function on $U$; and 
				\item[(ii)] $F(\mathbf{y}) - F(\mathbf{x})\leq
                DF(\mathbf{x})(\mathbf{y}-\mathbf{x})$ for all
                $\mathbf{x,y}\in U$; and
				\item[(iii)] $D^2F(\mathbf{x})$ is negative semidefinite
                for all $\mathbf{x}\in U$
			\end{enumerate}
			\item[(b)] The following three conditions are equivalent:
			\begin{enumerate}
				\item[(i)] $F$ is a convex function on $U$; and
				\item[(ii)] $F(\mathbf{y}) - F(\mathbf{x})\geq
                DF(\mathbf{x})(\mathbf{y}-\mathbf{x})$ for all
                $\mathbf{x,y}\in U$; and
				\item[(iii)] $D^2F(\mathbf{x})$ is positive semidefinite
                for all $\mathbf{x}\in U$
			\end{enumerate}
            \item[(c)] If $F$ is concave on $U$ and
            $DF(\mathbf{x^*}) = \mathbf{0}$ for some $\mathbf{x^*}\in U$,
            then $\mathbf{x^*}$ is a global max of $F$ on $U$.
            \item[(d)] If $F$ is convex on $U$ and $DF(\mathbf{x^*}) =
            \mathbf{0}$ for some $\mathbf{x^*}\in U$, then $\mathbf{x^*}$ is
            a global min of $F$ on $U$.
		\end{enumerate}	
	\end{block}
\end{frame}		
		

\begin{frame}
	\frametitle{Economic Applications: Discriminating Monopolist}
	\begin{itemize}
		\item A monopolist faces two distinct and separated markets (e.g., a
        domestic market and a foreign market), each with its own demand
        function
        \begin{itemize}
            \item supply: $Q_i$
            \item inverse demand function: $P_i = G_i(Q_i)$
            \item revenue: $Q_iG_i(Q_i)$
            \item production costs: $C(Q_1+Q_2)$
        \end{itemize}
		\item Profit: $F(Q_1,Q_2) = Q_1G_1(Q_1)+Q_2G_2(Q_2)-C(Q_1+Q_2)$
        \item Suppose that the firm will produce a positive amount for each
        market
		\item Problem: compute the maxima of the profit function $F$ in the
        interior of the positive quadrant
	\end{itemize}
\end{frame}


\begin{frame}
	\frametitle{Economic Applications: Discriminating Monopolist}
	\begin{itemize}
		\item Then we have
        $$\frac{d(Q_1G_1(Q_1))}{dQ_1} = \frac{d(Q_2G_2(Q_2))}{dQ_2} =
        C'(Q_1+Q_2)$$
		\item The marginal revenue in \emph{each} market = marginal cost of
        the total output
	\end{itemize}
\end{frame}


\begin{frame}
    \frametitle{Economic Applications: Discriminating Monopolist}
    \framesubtitle{Example}
    \begin{itemize}
        \item $G(Q_1) = 50 - 5Q_1$
        \item $G(Q_2) = 100 - 10Q_2$
        \item $C(Q) = 90 + 20Q$
        \item In order to maximize profits, how much should the monopolist
        produce for each market?
    \end{itemize}
\end{frame}


\begin{frame}
    \frametitle{Economic Applications: Discriminating Monopolist}
    \framesubtitle{Example}
    \begin{itemize}
		\item This discriminating monopolist’s profit function is 
		\begin{equation*}
            F(Q_1,Q_2) = Q_1(50-5Q_1)+Q_2(100-10Q_2) - (90+20(Q_1+Q_2))
		\end{equation*}
        \item The critical point of $F$ satisfy: 
        \begin{align*}
            \frac{\partial F}{\partial Q_1} &= 50-10Q_1-20 = 0 \\
            \frac{\partial F}{\partial Q_2} &= 100-20Q_2-20 = 0 
        \end{align*}
        Hence, $Q_1 = 3$ and $Q_2 = 4$
    \end{itemize}
\end{frame}


\begin{frame}
    \frametitle{Economic Applications: Discriminating Monopolist}
    \framesubtitle{Example}
	\begin{itemize}
		\item Now check the second order conditions: $F_{Q_1Q_1} = -10,
        F_{Q_2Q_2} = -20, F_{Q_1Q_2} = F_{Q_2Q_1} = 0$
		\item The first order leading principal minor of $D^2F(3,4)$ is -10
        and second order leading principal minor is 200.
		\item Therefore, F is a concave function and the point $(3, 4)$ is a
        maximizer
	\end{itemize}
\end{frame}

\end{document}
